\subsection{Composition products and connecting homomorphisms of torsion products}

PROPOSITION 7. --- a) Let

\[
0 \to \mathrm{P} ^ {\prime} \xrightarrow {f} \mathrm{P} \xrightarrow {g} \mathrm{P} ^ {\prime \prime} \to 0\tag{8}
\]

be an exact sequence of right A-modules, \(\theta\in\mathrm{Ext}_{\mathrm{A}}^{1}\left(\mathrm{P}^{\prime\prime},\mathrm{P}^{\prime}\right)\) the associated class, M a left A-module. The connecting homomorphism

\[
\delta_ {n} (\mathcal {E}, \mathrm{M}): \operatorname{Tor} _ {n} ^ {\mathrm{A}} \left(\mathrm{P} ^ {\prime \prime}, \mathrm{M}\right)\rightarrow \operatorname{Tor} _ {n - 1} ^ {\mathrm{A}} \left(\mathrm{P} ^ {\prime}, \mathrm{M}\right) \quad \text { is the map } \quad\gamma \mapsto \theta \circ \gamma .
\]

\begin{enumerate}
\def\labelenumi{\alph{enumi})}
\setcounter{enumi}{1}
\tightlist
\item
  Let
\end{enumerate}

\[
0 \rightarrow \mathrm{M} ^ {\prime} \rightarrow \mathrm{M} \rightarrow \mathrm{M} ^ {\prime \prime} \rightarrow 0\tag{$(\mathcal{E}_1)$}
\]

be an exact sequence of left A-modules, \(\theta_{1} \in \operatorname{Ext}_{\mathbf{A}}^{1}(\mathbf{M}^{\prime \prime}, \mathbf{M}^{\prime})\) the associated class, P a right A-module. The connecting homomorphism

\[
\delta_ {n} (\mathrm{P}, \mathcal {E} _ {1}): \operatorname{Tor} _ {n} ^ {\mathrm{A}} (\mathrm{P}, \mathrm{M} ^ {\prime \prime}) \rightarrow \operatorname{Tor} _ {n - 1} ^ {\mathrm{A}} (\mathrm{P}, \mathrm{M} ^ {\prime}) \quad \text { is the map } \quad\gamma \mapsto \theta_ {1} \circ \gamma .
\]

Let \(\gamma\in\mathrm{Tor}_{n}^{\mathrm{A}}(\mathrm{P}^{\prime\prime},\mathrm{M})\) the class of a cycle \(z^{\prime\prime}\in\dot{\mathrm{Z}}_{n}(\mathrm{L}(\mathrm{P}^{\prime\prime})\otimes_{\mathrm{A}}\mathrm{L}(\mathrm{M}))\) , and let

\[
\begin{array}{c} 0 \longrightarrow \mathrm{P} ^ {\prime} \xrightarrow {f} \mathrm{P} \xrightarrow {g} \mathrm{P} ^ {\prime \prime} \longrightarrow 0 \\ \uparrow_ {u _ {1}} \quad \uparrow_ {u _ {0}} \quad \uparrow_ {1} \\ \mathrm{L} _ {1} (\mathrm{P} ^ {\prime \prime}) \xrightarrow {d _ {1}} \mathrm{L} _ {0} (\mathrm{P} ^ {\prime \prime}) \xrightarrow {p _ {0} ^ {\prime \prime}} \mathrm{P} ^ {\prime \prime} \longrightarrow 0 \end{array}
\]

a commutative diagram. We will denote \(p' : \mathrm{L}(\mathrm{P}') \to \mathrm{P}'\) and \(p'' : \mathrm{L}(\mathrm{P}''') \to \mathrm{P}''\) the canonical morphisms of complexes. By definition, \(\delta(\gamma) \in \mathrm{Tor}_{n-1}^{\mathrm{A}}(\mathrm{P}', \mathrm{M})\) is obtained as follows: one chooses \(x \in \mathrm{P} \otimes \mathrm{L}_n(\mathrm{M})\) such that \((g \otimes 1)(x) = (p'' \otimes 1)(z'')\) and \(\delta(\gamma)\) is the class of cycles \(z' \in Z_{n-1}(\mathrm{L}(\mathrm{P}') \otimes \mathrm{L}(\mathrm{M}))\) such that

\[
(f \otimes 1) (p ^ {\prime} \otimes 1) (z ^ {\prime}) = (1 \otimes d _ {n}) (x).
\]

For \(0 \leqslant i \leqslant n\)  \(z_{i}^{\prime \prime}\) the component of \(z^{\prime \prime}\) to \(\mathrm{L}_{i}(\mathrm{P}^{\prime \prime}) \otimes \mathrm{L}_{n - i}(\mathrm{M})\) ; we have

\[
0 = \mathrm{D} z ^ {\prime \prime} = \sum_ {i} \left(d _ {i} \otimes 1 + (- 1) ^ {i} \otimes d _ {n - i}\right) (z _ {i} ^ {\prime \prime}),
\]

\(\mathrm{donc}(d_i \otimes 1)(z_i'') = (-1)^i \otimes d_{n-i+1}(z_{i-1}'')\) and in particular

\[
\left(d _ {1} \otimes 1\right) \left(z _ {1} ^ {\prime \prime}\right) = - 1 \otimes d _ {n} \left(z _ {0} ^ {\prime \prime}\right).
\]

Let us then choose \(x = (u_0 \otimes 1)(z_0'')\) : we indeed have

\[
(g \otimes 1) (x) = \left(p _ {0} ^ {\prime \prime} \otimes 1\right) z _ {0} ^ {\prime \prime} = \left(p ^ {\prime \prime} \otimes 1\right) \left(z ^ {\prime \prime}\right).
\]

Since

\[
\begin{array}{r l} (1 \otimes d _ {n}) (x) & = (u _ {0} \otimes 1) (1 \otimes d _ {n}) (z _ {0} ^ {\prime \prime}) = - (u _ {0} \otimes 1) (d _ {1} \otimes 1) (z _ {1} ^ {\prime \prime}) \\ & = - (f \otimes 1) (u _ {1} \otimes 1) (z _ {1} ^ {\prime \prime}), \end{array}
\]

it follows that \(\delta(\gamma)\) is the class of cycles \(z' \in Z_{n-1}(\mathbf{L}(\mathbf{P}') \otimes_{\mathbf{A}} \mathbf{L}(\mathbf{M}))\) such that \((p' \otimes 1)(z') = -(u_1 \otimes 1)(z_1'')\) . But, by definition, the class \(\theta\) corresponds under the isomorphism \(\operatorname{Ext}_A^1(\mathbf{P}', \mathbf{P}') \to \mathrm{H}^1(\operatorname{Homgr}_A(\mathbf{L}(\mathbf{P}''), \mathbf{P'}))\) to the class of the morphism \(f: \mathbf{L}(\mathbf{P}') (1) \to \mathbf{P}'\) defined by \(-u_1\) , and the product \(\theta \circ \gamma\) is the class of cycles

\[
\overline {{{z}}} ^ {\prime} \in Z _ {n - 1} (L (P ^ {\prime}) \otimes_ {A} L (M)) \quad \text { tels   que } \quad (p \otimes 1) (\overline {{{z}}} ^ {\prime}) = f (z ^ {\prime \prime}) = - (u _ {1} \otimes 1) (z _ {1} ^ {\prime \prime}),
\]

which completes the proof of \(a\) ). The assertion \(b\) ) follows from \(a\) ) by the commutation isomorphisms.

COROLLARY 1. --- Let \(0 \to P' \to P \to P'' \to 0\) be an exact sequence of right A-modules, \(0 \to M' \to M \to M'' \to 0\) an exact sequence of left A-modules. Then the homomorphisms composed of connecting homomorphisms

\[
\operatorname{Tor} _ {n} ^ {\mathbf {A}} \left(\mathrm{P} ^ {\prime \prime}, \mathrm{M} ^ {\prime \prime}\right)\rightarrow \operatorname{Tor} _ {n - 1} ^ {\mathbf {A}} \left(\mathrm{P} ^ {\prime \prime}, \mathrm{M} ^ {\prime}\right)\rightarrow \operatorname{Tor} _ {n - 2} ^ {\mathbf {A}} \left(\mathrm{P} ^ {\prime}, \mathrm{M} ^ {\prime}\right)
\]

and

\[
\operatorname{Tor} _ {n} ^ {\mathbf {A}} \left(\mathrm{P} ^ {\prime \prime}, \mathrm{M} ^ {\prime \prime}\right)\rightarrow \operatorname{Tor} _ {n - 1} ^ {\mathbf {A}} \left(\mathrm{P} ^ {\prime}, \mathrm{M} ^ {\prime \prime}\right)\rightarrow \operatorname{Tor} _ {n - 2} ^ {\mathbf {A}} \left(\mathrm{P} ^ {\prime}, \mathrm{M} ^ {\prime}\right)
\]

are opposite.

Indeed, if \(\theta\) and \(\theta_{1}\) are the classes associated with the given exact sequences, and if \(\gamma\in\mathrm{Tor}_{n}^{\mathbf{A}}(\mathbf{P}^{\prime\prime},\mathbf{M}^{\prime\prime})\) the images of \(\gamma\) are respectively \(\theta\circ(\theta_{1}\circ\gamma)\) and \(\theta_{1}\circ(\theta\circ\gamma)\) , hence are opposite by prop. 6.

Let us take up the notations of X, p.~127 and consider the sequence (S) of left A-modules and the connecting homomorphisms associated with the exact sequences (9).

\[
\operatorname{Tor} _ {m} ^ {\mathbf {A}} \left(\mathrm{P}, \mathbf {K} _ {i - 1}\right)\rightarrow \operatorname{Tor} _ {m - 1} ^ {\mathbf {A}} \left(\mathrm{P}, \mathbf {K} _ {i}\right);
\]

we deduce by composing the iterated connecting homomorphisms

\[
\partial_ {m} (\mathrm{P}, \mathscr {S}): \operatorname{Tor} _ {m} ^ {\mathrm{A}} (\mathrm{P}, \mathrm{M}) \rightarrow \operatorname{Tor} _ {m - n} ^ {\mathrm{A}} (\mathrm{P}, \mathrm{N}).
\]

Then by prop. 7 and prop. 3 of X, p.~118:

Corollary 2. --- If \(\theta \in \operatorname{Ext}_{\mathbf{A}}^{n}(\mathbf{M}, \mathbf{N})\) is the class associated with the exact sequence \((\mathcal{S})\) , we have \(\partial_{m}(\mathrm{P}, \mathcal{S})(\alpha) = \theta \circ \alpha\) for all \(\alpha \in \operatorname{Tor}_{m}^{\mathbf{A}}(\mathrm{P}, \mathbf{M})\) .

COROLLARY 3. --- If all the modules \(\mathbb{R}_i\) , \(i = 1, \ldots, n\) , are flat, the map \(\alpha \mapsto \theta \circ \alpha\) of \(\operatorname{Tor}_{m+n}^{\mathbf{A}}(\mathbf{P}, \mathbf{M})\) to \(\operatorname{Tor}_m^{\mathbf{A}}(\mathbf{P}, \mathbf{N})\) is bijective for every right A-module P and every integer \(m > 0\) .

This follows from Cor. 2 and the exact sequences.

\[
\operatorname{Tor} _ {m + n - i + 1} ^ {\mathbf {A}} \left(\mathrm{P}, \mathrm{R} _ {i}\right)\rightarrow \operatorname{Tor} _ {m + n - i + 1} ^ {\mathbf {A}} \left(\mathrm{P}, \mathrm{K} _ {i - 1}\right) \xrightarrow {\hat {c}} \operatorname{Tor} _ {m + n - i} ^ {\mathbf {A}} \left(\mathrm{P}, \mathrm{K} _ {i}\right)\rightarrow \operatorname{Tor} _ {m + n - i} ^ {\mathbf {A}} \left(\mathrm{P}, \mathrm{R} _ {i}\right)
\]

where the extreme terms are zero by hypothesis.

Likewise, if

\[
0 \rightarrow \mathrm{Q} \rightarrow \mathrm{S} _ {n} \rightarrow \mathrm{S} _ {n - 1} \rightarrow \dots \rightarrow \mathrm{S} _ {1} \rightarrow \mathrm{P} \rightarrow 0\tag{$(\mathcal{S}_1)$}
\]

is an exact sequence of right A-modules, and M is a left A-module, one defines iterated connecting homomorphisms

\[
\partial^ {m} (\mathcal {S} _ {1}, \mathrm{M}): \operatorname{Tor} _ {m} ^ {\mathrm{A}} (\mathrm{P}, \mathrm{M}) \rightarrow \operatorname{Tor} _ {m - n} ^ {\mathrm{A}} (\mathrm{Q}, \mathrm{M})
\]

and we have:

COROLLARY 4. --- If \(\theta_{1} \in \operatorname{Ext}_{\mathbf{A}}^{n}(\mathbf{P}, \mathbf{Q})\) is the class associated with the exact sequence \((\mathcal{S}_{1})\) , we have \(\partial^{m}(\mathcal{S}_{1}, \mathbf{M})(\alpha) = \theta_{1} \circ \alpha\) for all \(\alpha \in \operatorname{Tor}_{m}^{\mathbf{A}}(\mathbf{P}, \mathbf{M})\) .

